\section{Reproducibility and Deliverable Links}
\label{app:reproducibility}

\subsection{Artifact provenance and run manifest}
% Check recorded revisions against the source result JSONs in results/a1/.
\begin{longtable}{p{0.14\linewidth}p{0.46\linewidth}p{0.30\linewidth}}
\caption{Run provenance snapshot. Paths are relative to the repository root.
The configuration file for each model is under
\codepath{configs/image_classification/fashion_mnist/}.}
\label{tab:manifest}\\
\toprule Model & Experiment / source records & Recorded source revision\\\midrule
\endfirsthead
\toprule Model & Experiment / source records & Recorded source revision\\\midrule\endhead
Linear & \codepath{fashion_mnist_linear_seed42}\newline \codepath{results/a1/fashion_mnist_linear_seed42/} & \codepath{41c2424436b4} + dirty tree \\
MLP & \codepath{fashion_mnist_mlp_seed42}\newline \codepath{results/a1/fashion_mnist_mlp_seed42/} & \codepath{41c2424436b4} + dirty tree \\
CNN & \codepath{fashion_mnist_cnn_spatial_seed42}\newline \codepath{results/a1/fashion_mnist_cnn_spatial_seed42/} & \codepath{e6de6df2fff9} + dirty tree \\
GRU & \codepath{fashion_mnist_gru_seed42}\newline \codepath{results/a1/fashion_mnist_gru_seed42/} & \codepath{41c2424436b4} + dirty tree \\
LSTM & \codepath{fashion_mnist_lstm_seed42}\newline \codepath{results/a1/fashion_mnist_lstm_seed42/} & \codepath{c89aacfe7721} + dirty tree \\
Transformer & \codepath{fashion_mnist_transformer_seed42}\newline \codepath{results/a1/fashion_mnist_transformer_seed42/} & \codepath{41c2424436b4} + dirty tree \\
\bottomrule
\end{longtable}

\reporttodo{Associate every result with its exact configuration, dataset split
indices or reproducible split recipe, checkpoint, log/experiment ID and source
commit/tag. A recorded revision ending in -dirty is incomplete provenance:
preserve the corresponding working-tree patch or reproduce from a clean tagged
revision, and document which evidence supports the final claim.}
The JSON snapshots under \codepath{Evidence/} preserve source records used for
these tables. The numerical snapshot is not an additional experiment.

\subsection{Installation, training and evaluation}
Execute from the repository root, not this LaTeX directory. The following shows
one model; repeat with each verified model configuration.
\begin{lstlisting}[basicstyle=\ttfamily\scriptsize,breaklines=true,frame=single]
uv sync --dev
uv run python -m scripts.eda.run_fashion_mnist
uv run python -m scripts.training.train_image_classification --data fashion_mnist --config configs/image_classification/fashion_mnist/linear.yaml
uv run python -m scripts.evaluation.evaluate --data fashion_mnist --config configs/image_classification/fashion_mnist/linear.yaml
uv run python -m scripts.evaluation.compare --results-dir results/a1
\end{lstlisting}
\reporttodo{Verify these commands in the final environment. Explain dataset download
and integrity checking, checkpoint reconstruction, exact environment/dependency
versions and hardware. Ensure the main comparison contains only the intended
Fashion-MNIST runs; document filtering if the results directory also contains
MNIST or CIFAR-10 experiments. Add checkpoint URLs or reconstruction instructions.}

\subsection{Submission links and individual contributions}
\begin{table}[htbp]
\centering\small
\caption{Submission artifact index. Replace every pending entry with a checked link.}
\begin{tabularx}{\linewidth}{lY}
\toprule
Artifact & Location\\\midrule
Source repository & \href{\RepoURL}{Source repository}\\
Project landing page & \href{\ProjectURL}{Project landing page}\\
Assignment 1 page & \url{https://vietanh1802.github.io/CO3133-Deep-Learning-261-Placeholder/web/assignment1.html}\\
Report / slides & \pending\\
YouTube presentation & \pending; Public or Unlisted; all members participate\\
Checkpoints & \pending; or give verified reconstruction instructions\\
AI disclosure & Appendix~\ref{app:ai}, repository \codepath{AI_USAGE.md}, and both web pages\\
Final PDF on LMS & \codepath{Placeholder_A1_Report.pdf}; representative submits\\
Updated A1 deadline & 28 October 2026, 23:59 GMT+7 (group update)\\
\bottomrule
\end{tabularx}
\end{table}
\reporttodo{Record each member's actual contribution and final verification
responsibilities. Confirm the registered group spelling and all external links.
Slides, video and website remain separate required deliverables.}

\clearpage
\subsection{Per-model learning curves and confusion matrices}
\label{app:evidence}
The following figures are copies of local Fashion-MNIST run artifacts. Each model
has its own training/validation curve and test confusion matrix.
% Included by the Reproducibility appendix, not a fourth main report part.
\subsubsection{LINEAR}
\reportfigure{Images/results/linear/learning_curves.png}{LINEAR: saved training and validation curves.}{fig:linear-curves}
\reportfigure{Images/results/linear/test_confusion_matrix.png}{LINEAR: saved test confusion matrix. Verify whether values are counts or normalized rates.}{fig:linear-confusion}
\clearpage

\subsubsection{MLP}
\reportfigure{Images/results/mlp/learning_curves.png}{MLP: saved training and validation curves.}{fig:mlp-curves}
\reportfigure{Images/results/mlp/test_confusion_matrix.png}{MLP: saved test confusion matrix. Verify whether values are counts or normalized rates.}{fig:mlp-confusion}
\clearpage

\subsubsection{CNN}
\reportfigure{Images/results/cnn/learning_curves.png}{CNN: saved training and validation curves.}{fig:cnn-curves}
\reportfigure{Images/results/cnn/test_confusion_matrix.png}{CNN: saved test confusion matrix. Verify whether values are counts or normalized rates.}{fig:cnn-confusion}
\clearpage

\subsubsection{GRU}
\reportfigure{Images/results/gru/learning_curves.png}{GRU: saved training and validation curves.}{fig:gru-curves}
\reportfigure{Images/results/gru/test_confusion_matrix.png}{GRU: saved test confusion matrix. Verify whether values are counts or normalized rates.}{fig:gru-confusion}
\clearpage

\subsubsection{LSTM}
\reportfigure{Images/results/lstm/learning_curves.png}{LSTM: saved training and validation curves.}{fig:lstm-curves}
\reportfigure{Images/results/lstm/test_confusion_matrix.png}{LSTM: saved test confusion matrix. Verify whether values are counts or normalized rates.}{fig:lstm-confusion}
\clearpage

\subsubsection{TRANSFORMER}
\reportfigure{Images/results/transformer/learning_curves.png}{TRANSFORMER: saved training and validation curves.}{fig:transformer-curves}
\reportfigure{Images/results/transformer/test_confusion_matrix.png}{TRANSFORMER: saved test confusion matrix. Verify whether values are counts or normalized rates.}{fig:transformer-confusion}
\clearpage

